Ntc\documentclass[11pt]{article}

\usepackage[margin=0.9in]{geometry}
\usepackage{amsmath,amssymb}
\usepackage{enumitem}
\usepackage[hidelinks]{hyperref}
\usepackage{titlesec}

\linespread{0.96}

\setlength{\parskip}{3pt}
\setlength{\parindent}{0pt}

\setlist[itemize]{itemsep=1pt,topsep=1pt}
\setlist[enumerate]{itemsep=1pt,topsep=1pt}

\titlespacing{\section}{0pt}{6pt}{3pt}

\begin{document}
	
	\begin{center}
		
		{\Large\bfseries Mini-Research Proposal}
		
		\vspace{0.2cm}
		
		{\large\bfseries Empirical Diagnostics of Diffusion and Jump Dynamics}
		
		{\large\bfseries in Equity Markets using Stochastic Differential Equations}
		
		\vspace{0.4cm}
		
		Course: \textbf{Stochastic Differential Equations (MATH565C)}
		
		% \vspace{0.2cm}
		
		Prepared by: \textbf{Dmitri Bolt}
		
		Instructor: \textbf{Prof. Ibrahim Fatkullin (University of Arizona)}
		
		Date: March 5, 2026
		
	\end{center}
	
	\vspace{0.3cm}
	
	\section*{Part I: Overview}
	
	This mini-research uses real equity price data (via the \texttt{yfinance} Python library) to examine how well simple continuous-time stochastic differential equation (SDE) models describe empirical log-returns.
	
	The baseline hypothesis is that, after the log transformation, a large portion of the random variation can be approximated by an Itô diffusion driven by a Wiener process \(W_t\), while rare large moves may require an additional jump component.
	
	The project develops hands-on familiarity with core tools from Øksendal’s textbook (Brownian motion / Wiener process, Itô integrals, Itô’s formula, quadratic variation) by translating them into empirical diagnostics.
	
	\textbf{Key questions}
	
	\begin{enumerate}
		\item Are log-return increments approximately Brownian at a given time scale (variance scaling with \(\Delta t\), weak serial dependence, approximate Gaussianity)?
		\item Do quadratic-variation diagnostics indicate mostly continuous diffusion behavior, or do a few large increments dominate (jump-like behavior)?
		\item Which simple SDE baseline best captures a market factor and a representative subset of stocks: geometric Brownian motion, Ornstein–Uhlenbeck dynamics on a de-trended factor, or diffusion with sparse jumps?
	\end{enumerate}
	
	\textbf{Approach}
	
	Download prices → compute log-returns → optionally extract a market factor (PCA) → run diffusion and jump diagnostics → fit simple SDE models → compare results.
	
	\textbf{Deliverables}
	
	A reproducible Python notebook, a short report (5–10 pages), and an explicit mapping between empirical steps and the corresponding concepts in Øksendal.
	
	\textbf{Rationale}
	
	The scope is intentionally small but exercises the mathematical core of stochastic differential equations while connecting empirical data analysis with diffusion theory.
	
	Each fitted diffusion model defines a generator
	\[
	L f(x)=a(x)f'(x)+\tfrac12 b(x)^2 f''(x),
	\]
	which determines the associated Fokker–Planck equation governing the evolution of the probability density.
	
	\[
	\text{SDE dynamics}
	\;\longrightarrow\;
	\text{diffusion generator}
	\;\longrightarrow\;
	\text{Fokker--Planck PDE}
	\]
	
	This viewpoint highlights the link between stochastic diffusion models and their PDE descriptions.
	
	\newpage
	
	\section*{Part II: Detailed Plan}
	
	\textbf{1. Study setup and data scope}
	
	Data consist of daily adjusted close prices obtained from Yahoo Finance using the \texttt{yfinance} Python library, covering approximately 20–200 liquid US equities together with an index proxy (e.g. SPY), which keeps the analysis computationally manageable while still representing market dynamics.
	
	Variables used in the study:
	
	\[
	S_t,\qquad X_t=\ln S_t,\qquad \Delta X_k=X_{t_{k+1}}-X_{t_k}.
	\]
	
	\textbf{2. Research hypotheses}
	
	Diffusion hypothesis
	
	\[
	dX_t=a\,dt+b\,dW_t
	\]
	
	implying
	
	\[
	\mathbb{E}[\Delta X]\approx a\Delta t,
	\qquad
	\mathrm{Var}(\Delta X)\approx b^2\Delta t.
	\]
	
	Jump hypothesis
	
	Heavy tails or large outliers may correspond to a sparse jump component visible through large contributions to realized quadratic variation.
	
	Mean-reversion hypothesis
	
	A de-trended market factor may follow Ornstein–Uhlenbeck dynamics
	
	\[
	dX_t=\kappa(\theta-X_t)dt+\sigma dW_t.
	\]
	
	\textbf{3. Mathematical toolkit from Øksendal}
	
	Primary reference:
	
	Øksendal, \emph{Stochastic Differential Equations: An Introduction with Applications}, 5th ed.
	
	Concepts used in the project:
	
	\begin{itemize}
		\item Wiener process (Brownian motion) — Section 2.2
		\item Itô integral — Section 3.1
		\item Itô processes and Itô’s formula — Chapter 4
		\item Quadratic variation \((X,X)_t\) — Chapter 4
		\item Geometric Brownian motion and log transform — Chapter 5
		\item Ornstein–Uhlenbeck process — Chapter 5 exercises
		\item Optional Stratonovich vs Itô conversion — Section 3.3
	\end{itemize}
	
	\textbf{4. Empirical methodology}
	
	Preprocessing
	
	Compute log-returns, treat missing values, and generate basic descriptive statistics.
	
	Diffusion diagnostics
	
	\begin{itemize}
		\item variance scaling across time aggregation
		\item serial dependence tests
		\item Gaussianity diagnostics
	\end{itemize}
	
	Quadratic variation
	
	\[
	QV(T)=\sum_k (\Delta X_k)^2
	\]
	
	used to identify whether a few extreme increments dominate the realized variance.
	
	Model fitting
	
	Estimate parameters for GBM (prices) and OU dynamics (de-trended factor).
	
	Optional jump extension
	
	Add a sparse jump component and compare against diffusion-only baselines.
	
	\textbf{5. Skills developed}
	
	\begin{enumerate}
		\item Understanding Wiener increments as drivers of SDE noise
		\item Practical use of Itô integrals and Itô’s formula
		\item Quadratic variation as a bridge between SDE theory and realized variance
		\item Translating continuous-time models to discrete-time data
		\item Reproducible computational workflow for SDE modeling
	\end{enumerate}
	
	\textbf{6. Proposed timeline}
	
	March 5 – March 15
	
	Dataset construction, preprocessing pipeline, initial diagnostics.
	
	March 16 – March 25
	
	Diffusion diagnostics (variance scaling, serial dependence, Gaussianity).
	
	March 26 – April 5
	
	Quadratic variation analysis and identification of large increments.
	
	April 6 – April 15
	
	Calibration of GBM and OU models and simulation checks.
	
	April 16 – April 25
	
	Optional jump extension and model comparison.
	
	April 26 – May 3
	
	Final report preparation and reproducible notebook cleanup.
	
	May 4
	
	Project presentation.
	
	\vspace{0.3cm}
	
	\begin{thebibliography}{1}
		
		\bibitem{oksendal}
		B. Øksendal,
		\emph{Stochastic Differential Equations: An Introduction with Applications},
		5th ed., Springer, 2003.
		
	\end{thebibliography}
	
\end{document}